\subsection{THETA3-A11: an explicit host extension and actual finite-jet realization}
\label{A11}

The finite CM host condition can be achieved constructively.  The
construction below retains the actual CM constituents, produces a
principal polarization and the archive's primitive symmetric line,
and then realizes a prescribed finite family of first jets.  Its
output is a generic \texttt{Host}/\texttt{Datum}/\texttt{Obj}.  An
identification with a particular \texttt{StdPair}, oscillator kernel,
or global iterated theta chain is additional data.

\paragraph{The exact source conditions.}
Put
\[
 \Psi=\{\tau:E\longrightarrow\C:\tau(\mathrm{rt})=\mathrm{rt}_0\}.
\]
The recovered source
\texttt{ReciprocityTheory/ReductiveGroups/ArithmeticGeometry/}
\texttt{BirationalGeometry/ToroidalCompactification/}
\texttt{TwoWeightEqOneAddRoot.lean},
\texttt{Host}, \texttt{Datum}, \texttt{Level}, \texttt{SAB},
and \texttt{Obj}, requires the following objects.  The field host
contains an integral square root of $-d$, where $d$ is a prime
congruent to $7$ modulo $16$, and $w_0=(1+\mathrm{rt})/2$ is
integral.  The extension over the maximal real subfield is
unramified at every finite prime.  A datum is an actual finite
projective $\OO_E$-lattice with a unimodular alternating integral
form, with adjoint involution complex conjugation, and a compatible
positive complex structure.  Its Lie multiplicities are the same
integer $p$ at every embedding in $\Psi$ and the same integer $q$
at the conjugate embeddings.  There is no condition $p\geq q$ or
parity condition in this definition.  The exact multiplicity
function is also written in \texttt{FlagBasis.lean},
\texttt{Datum.mult}.  The additional line-bundle identity is
\[
       L\simeq(\mathrm{id},[w_0])^*P_L,\qquad
       P_L=m^*L\otimes\operatorname{pr}_1^*L^{-1}
                    \otimes\operatorname{pr}_2^*L^{-1},
 \tag{A11a}
\]
with the origin rigidification understood.  This sign is the
literal two-variable \texttt{cubeLB} in
\texttt{ReciprocityTheory/ReductiveGroups/ArithmeticGeometry/}
\texttt{SemiabelianSchemes/CubicalStructures/}
\texttt{CubeSymmetricSwapFirst.lean}.

\paragraph{A universal field host of degree at most four.}
Let $E$ be any CM field.  Define
\[
 K=\Q(\sqrt{-7}),\qquad F_H=\Q(\sqrt{42}),\qquad
 H=K F_H,\qquad E'=EH,\qquad F'=(E')^+.
 \tag{A11b}
\]
Then
\[
 [E':E]\leq4,\qquad
 E'=F'K,\qquad
 \OO_{E'}\text{ is finite \'etale over }\OO_{F'}.
 \tag{A11c}
\]
Here $\mathrm{rt}=\sqrt{-7}$ and $w_0=(1+\sqrt{-7})/2$ give
all the field-generator conditions with $d=7$.

We prove the finite-prime assertion directly.  The quadratic field
$K$ has discriminant $-7$, so its integer ring is finite
\'etale away from $7$.  Base change to $F_H$ therefore gives
an unramified quadratic algebra away from $7$.  At $7$, the element
$-6=42/(-7)$ is a square in $\Q_7$: apply Hensel's lemma to
$X^2+6$ at $X=1$, whose derivative is a unit modulo $7$.
Consequently
\[
 \Q_7(\sqrt{42})=\Q_7(\sqrt{-7}),
\]
and $H/F_H$ is split at the prime above $7$.  It is thus unramified
at every finite prime.  The discriminant check is
$(-7)(-24)(168)=168^2$ for the three quadratic subfields, in
agreement with this local proof.

A compositum of CM fields is CM, so $E'$ is CM.  Its maximal real
subfield contains $F_H$, and $K$ is not contained in that real
field.  Since $[E':F']=2$, one has $E'=F'K$.  The finite
\'etale $\OO_{F'}$-algebra
$\OO_H\otimes_{\OO_{F_H}}\OO_{F'}$ has generic algebra $E'$.
It is normal, because finite \'etale algebras over a Dedekind
domain are normal.  It therefore equals the integral closure
$\OO_{E'}$.  This proves (A11c), including the integral assertion
required by \texttt{Host.unram}.  Finite \'etale base change
is Stacks, Lemma 29.37.4, tag \texttt{02GO}.  The CM compositum
and Galois-closure facts are Milne, \emph{Complex Multiplication}
(2020), Proposition 1.4 and Corollary 1.5, printed p.~10.

If $E/\Q$ is Galois, as in the original 032 setup, $E'/\Q$ is
already Galois.  In general one may replace $E'$ by its CM Galois
closure.  This closure still contains $H$, so precisely the same
base-change argument proves finite unramifiedness over its own
maximal real subfield.  The bound $[E':E]\leq4$ refers to (A11b),
before this optional Galois closure.

The same construction is available for any prescribed prime
$d\equiv7\pmod {16}$.  Replace $H$ by
$H_d=\Q(\sqrt{-d},\sqrt{d(d-1)})$ and its real subfield by
$\Q(\sqrt{d(d-1)})$.  Away from $d$ the quadratic field of
discriminant $-d$ is unramified, and at $d$ the ratio
$d(d-1)/(-d)=1-d$ is a square in $\Q_d$ by Hensel's lemma
at $1$.  The same local and integral base-change proof applies.
This variant permits a prescribed size condition on $d$
once such a prime has been chosen; it does not itself supply
the additional arithmetic parameters of the global archive.

\paragraph{Scalar extension of the actual CM abelian variety.}
Let the positive-dimensional $A/\overline\Q$ have an action of $R=\OO_E$ and an
$E$-compatible polarization.  Let
$s=[E':E]$, $S=\OO_{E'}$, and $S^*=\Hom_R(S,R)$.  Define
\[
       A'=\Hom_R(S^*,A),
       \qquad H_1(A',\Q)=H_1(A,\Q)\otimes_E E'.
 \tag{A11d}
\]
The first expression denotes a represented group functor.
Since $S$ is finite projective over the Dedekind domain $R$,
$S^*$ is a direct summand of a finite free module.  Applying
$\Hom_R(-,A)$ to a splitting realizes $A'$ as an actual direct
factor of a power of $A$.  It is therefore an abelian variety.
Multiplication on $S$ gives its $S$-action.  The same idempotent
construction on homology proves the second equality, and gives
\[
       \dim A'=s\dim A,\qquad A'\sim A^s.
 \tag{A11e}
\]
Here the isogeny is $E$-linear; it is not an assertion that $A$
already carried the larger field action.  This is the projective
Serre construction in Milne, \emph{Complex Multiplication},
Proposition 7.33, Remark 7.34, and Proposition 7.35, printed
pp.~62--63.  The proof just given also explains its compatibility
with abelian schemes and base change.

There is even an $R$-linear direct-factor inclusion of $A$ in
$A'$.  Indeed, $S/R$ is torsion free: if $0\ne a\in R$ and
$ax\in R$ for $x\in S$, then $x\in E$ and its integrality gives
$x\in R$.  Thus $S/R$ is projective and $R\cdot1\subset S$
splits.  Tensoring this splitting with $A$ gives the inclusion.
It is generally not $E'$-linear.  For an embedding $\tau'$ of
$E'$ extending $\tau$ of $E$, the induced Lie decomposition is
\[
       \operatorname{Lie}(A')_{\tau'}
                       =\operatorname{Lie}(A)_\tau.
 \tag{A11f}
\]
More precisely, project the preceding inclusion to the
$\tau'$-eigenspace; its restriction to the $\tau$-space is the
displayed isomorphism.  All these maps and eigenspaces are
algebraic over a finite number field.  Integrality of an
eigenspace projection is only asserted after its denominators
are inverted, or at an unramified split coefficient base where
the corresponding idempotent is integral.

For completeness, the extended action has a compatible
polarization.  If $\psi$ is the original rational alternating
Riemann form, trace nondegeneracy gives a unique $E$-valued form
$h$ with
\[
       \Tr_{E/\Q}(a h(x,y))=\psi(ax,y),\qquad a\in E.
\]
The adjoint identity for $\psi$ implies that $h$ is linear in
the first variable, conjugate linear in the second, and
$h(y,x)=-\overline{h(x,y)}$.  Extend $h$ to $E'$ and use
$\Tr_{E'/\Q}h$.  At every complex embedding this is the same
positive Riemann form as the original one, on each of the
additional summands.  It satisfies the Riemann relations.
Multiplication by a positive integer makes it integral on
$H_1(A',\Z)$ and supplies the asserted polarization.

\paragraph{Minimal uniform multiplicity padding.}
For this paragraph write $E$ for the enlarged field and
$r=\dim_E H_1(A',\Q)$.  Put
$m_\tau=\dim_{\overline\Q}\operatorname{Lie}(A')_\tau$.
The Hodge decomposition and conjugation give
$m_\tau+m_{c\tau}=r$.  Define
\[
 P=\max_{\tau\in\Psi}m_\tau,\qquad
 m=\min_{\tau\in\Psi}m_\tau,\qquad
 Q=r-m,\qquad R=P-m.
 \tag{A11g}
\]
There are $R$ rank-one CM abelian varieties $B_1,\ldots,B_R$
such that
\[
        B=A'\times\prod_{j=1}^R B_j
 \tag{A11h}
\]
has multiplicities $(P,Q)$ at every pair $(\tau,c\tau)$ with
$\tau\in\Psi$.  To construct $B_j$, choose its CM type to
contain $\tau\in\Psi$ exactly when $j\leq P-m_\tau$, and
otherwise to contain $c\tau$.  This selects exactly one
embedding from every conjugate pair.  The added multiplicities
at that pair are $P-m_\tau$ and $m_\tau-m$, respectively,
which proves the assertion.

Every such type is realized by a polarized rank-one CM torus:
use an $\OO_E$-lattice in $E$, choose a purely imaginary
element with the signs prescribed by the type, and take the
trace alternating form.  Weak approximation in the maximal
real subfield supplies the signs, and a positive integer clears
its denominators.  The Riemann relations give an abelian
variety, and CM descent gives a model over $\overline\Q$
(Milne, Corollary 7.10, printed p.~55).  Weak approximation is
Milne, \emph{Algebraic Number Theory} (2020), Theorem 7.20,
printed p.~114.  These are actual added constituents.  No
isogeny between different full CM types is inferred from their
local multiplicity counts.

The number $R$ is minimal for padding by rank-one CM factors.
Indeed, any uniform target $(P',Q')$ must satisfy
$P'\geq P$ and $Q'\geq r-m$.  The number of added factors is
$P'+Q'-r$, hence is at least $P-m=R$.  In particular $R\leq r$.

\paragraph{An explicit compatible principal polarization.}
Let $\lambda:B\to B^\vee$ be the product compatible
polarization.  Choose $N\geq1$ annihilating $\ker\lambda$;
then $\lambda'=N\lambda^{-1}:B^\vee\to B$ is an actual
homomorphism.  Choose integers $a,b,c,d$ whose squares sum to
$N-1$, and form
\[
 T=\begin{pmatrix}
 a&b&c&d\\-b&a&-d&c\\-c&d&a&-b\\-d&-c&b&a
 \end{pmatrix},\qquad
 TT^t=T^tT=(N-1)I_4.
\]
Existence of the four integers is Lagrange's four-square
theorem.  Give $B^\vee$ the conjugate-dual $E$-action
$u\mapsto\iota_B(\bar u)^\vee$, and put
\[
 X=B^4\times(B^\vee)^4,\qquad
 \Lambda=
 \begin{pmatrix}\lambda I_4&T^t\\T&\lambda'I_4\end{pmatrix}
       :X\longrightarrow X^\vee.
 \tag{A11i}
\]
This is a principal polarization with Rosati involution $c$
on $E$.  We verify every assertion.  The matrix is symmetric,
and its inverse is the actual homomorphism
\[
 \Lambda^{-1}=
 \begin{pmatrix}\lambda'I_4&-T^t\\-T&\lambda I_4\end{pmatrix},
 \tag{A11j}
\]
as multiplication gives $N-(N-1)=1$ on the diagonal and zero
off the diagonal.  Positivity follows from the positive first
block and its Schur complement
\[
       \lambda'I_4-T(\lambda^{-1}I_4)T^t
                         =\lambda^{-1}I_4>0.
\]
Thus the symmetric map is a polarization, and (A11j) makes it
principal.  The equation
$\iota_B(u)^\vee\lambda=\lambda\iota_B(\bar u)$,
the analogous equation for $\lambda'$, and the integral scalar
entries of $T$ verify block by block that
\[
       \iota_X(u)^\vee\Lambda=\Lambda\iota_X(\bar u).
\]
Finally, $\lambda$ is $E$-linear from $B$ to the
conjugate-dual $B^\vee$.  The latter consequently has the same
full $E$-Hodge constituents and the same Lie multiplicities.
The uniform multiplicities of $X$ are $(8P,8Q)$.

Let $P_\Lambda$ be the normalized Poincar\'e biextension on
$X\times X$ classified by $\Lambda$, and define
\[
                 L=(\mathrm{id},[w_0])^*P_\Lambda.
 \tag{A11k}
\]
The polarization of this rigidified line is
\[
  \phi_L=\Lambda[w_0]+[w_0]^\vee\Lambda
          =\Lambda[w_0+\bar w_0]=\Lambda.
\]
Hence $L$ is ample and principal.  Pullback by $[-1]$ on both
arguments of a biextension leaves it unchanged, so $L$ is
symmetric.  Its normalized biextension $P_L$ is $P_\Lambda$,
which proves (A11a) as an isomorphism of actual rigidified
lines.  The theorem of the cube supplies its cubical structure,
and the adjoint relation follows from the displayed Rosati
identity.  These are the Poincar\'e/line-bundle constructions of
SGA~7 I, VII, Corollary 2.9.4 and Example 2.9.5, printed
pp.~174--175, and VIII, 3.1--3.2, printed pp.~254--255.

Take $\mathcal L=H_1(X,\Z)$ and its principal alternating
Riemann form, with the convention $EF(Jx,x)>0$ used in the
source.  The lattice is finite and torsion free over the
Dedekind domain $\OO_E$, hence projective.  Principalness gives
unimodularity.  The Rosati equation gives the required adjoint
identity.  Multiplication by $i$ gives $J$.  If
\[
 t=\sum_{\tau\in\Psi}\bigl(8P\tau(u)+8Q\overline{\tau(u)}\bigr),
\]
then the real traces are $2\operatorname{Re}t$ and
$-2\operatorname{Im}t$ for $u$ and $Ju$, exactly as in
\texttt{Datum.dom}.  This constructs the actual datum.

The dimension cost is explicit.  If $g=\dim A$ before the
field extension, then
\[
       \dim X=8s\frac{r+R}{r}\,g\leq64g,
 \tag{A11l}
\]
using $s\leq4$ and $R\leq r$.  An optional Galois closure of
a non-Galois original field replaces the bound on $s$ by its
actual extension degree.  The generic unitary moduli dimension
for the constructed datum is
$[E:\Q](8P)(8Q)/2$; this is the exact \texttt{Datum.dim}
in \texttt{PushInvertible.lean}.

\paragraph{Levels and an actual complex/algebraic object.}
There are arbitrarily large admissible odd levels avoiding any
one prescribed prime $p$.  Choose $\ell=11$ or $23$, with
$\ell\ne p$, and a prime $\mathfrak l$ of $K$ above $\ell$.
Both primes split in $K$: $X^2-X+2$ has roots $5,7$ modulo
$11$ and $10,14$ modulo $23$.  For $t\geq1$, put
\[
       M=\ell^t,\qquad
       \mathfrak M=\mathfrak l^t\OO_E,\qquad A_{\rm aux}=1.
 \tag{A11m}
\]
Then $\mathfrak M\overline{\mathfrak M}=(M)$ and
$\mathfrak M+\overline{\mathfrak M}=\OO_E$.  Also
$\mathfrak M$ is generated by its elements in
$\Z[w_0]=\OO_K$, proving precisely \texttt{Level.fromK}.
All integers are invertible over the base $\overline\Q$,
and $\rho=\sqrt{-7}$ and $\zeta=1$ supply \texttt{Base}.

For clarity, the remaining level data can be constructed
without a moduli existence assertion.  In
$X(\C)=\operatorname{Lie}(X)/\mathcal L$, set
$\mathrm{ev}_n(v)=v/n\pmod{\mathcal L}$, with a compatible
choice of roots of unity matching the Riemann-form convention.
These maps give \texttt{TateIso}: their kernel is
$n\mathcal L$, their image is all $n$-torsion, they commute
with the $\OO_E$-action, and they have the required tower and
Weil-pairing identities.  The Chinese remainder decomposition
\[
 \mathcal L/M\mathcal L
       =(\mathfrak M\mathcal L/M\mathcal L)
        \oplus(\overline{\mathfrak M}\mathcal L/M\mathcal L)
\]
defines $\eta$ by projecting to the second summand and then
applying $\mathrm{ev}_M$.  It kills
$\mathfrak M\mathcal L$ and equals $\mathrm{ev}_M$ on
$\overline{\mathfrak M}\mathcal L$, as required.  Take
$\eta_A=0$ when $A_{\rm aux}=1$.  The finite level points are
algebraic, so the object and its level structure are defined
over a finite extension of a number field.  The auxiliary
compatible torsion identifications for all $n$ are required
only over algebraically closed geometric fibres; their
simultaneous descent to that number field is not asserted.
Here is explicitly why the determinant condition survives
every coefficient embedding.  Enlarge the number field of
definition to contain $\rho=\sqrt{-7}$ and all the embedding
values of $E$.  For a field embedding $\beta$ into an
algebraically closed field of characteristic zero, extend
$\beta$ to an embedding of $\overline\Q$.  A Lie eigenspace
labelled by $\tau$ then becomes the eigenspace labelled by
$\beta\circ\tau$.  Since
\[
 (\beta\circ\tau)(\mathrm{rt})=\beta(\rho)
                  \quad\Longleftrightarrow\quad
                  \tau(\mathrm{rt})=\rho,
\]
its multiplicity is still $8P$ in this case and $8Q$
otherwise.  This is exactly \texttt{DetCond} with the
transported root $\beta(\rho)$; the root is not held fixed
while conjugating the object.  Transport every auxiliary
$\mathrm{ev}_n$ and its compatible root of unity through
the same extension of $\beta$.  The integral Riemann-form
entries, tower relations, and algebraic Weil pairings are
preserved.  The finite level maps already descend to the
number field, so the transported identifications agree with
their prescribed base changes.  This proves the entire
\texttt{geom} condition.  Principalness gives
\texttt{Out.KTriv}: a point whose translation preserves $L$
lies in the trivial kernel of $\phi_L$.  Thus the constructed
abelian variety, line, action, and level form an actual
\texttt{Obj}, over $\overline\Q$ and after finite descent.

\paragraph{Preserving a finite family of first-jet values.}
Changing the polarization in (A11i) changes its theta line.
The following precise interpolation argument accounts for
that change.  Let $U$ be the finite-dimensional algebraic input
space containing the finite class-sum inputs under consideration.
Transport each prescribed algebraic tangent direction through
(A11f), the product inclusion in (A11h), and one of the direct
factors of $X$.  This gives the selected tangent subspace
$V_0\subset\operatorname{Lie}(X)$.  Retain the algebraic
values and first derivatives of the old sections on $V_0$.
Choose an algebraic complement of $V_0$ and extend the desired
linear functionals by zero on that complement.

Let $Z_1$ be the first infinitesimal neighborhood of zero,
defined by the square of its ideal.  A normalized symmetric
origin frame identifies
\[
       H^0(Z_1,L^{M^2}|_{Z_1})
                    =\overline\Q\oplus\operatorname{Lie}(X)^*.
\]
The prescription is therefore one linear map from $U$ to this
finite-dimensional jet space.  For all sufficiently large
$M$ in the sequence (A11m), the map
\[
 H^0(X,L^{M^2})\longrightarrow H^0(Z_1,L^{M^2}|_{Z_1})
 \tag{A11n}
\]
is surjective.  Indeed, tensor the ideal sequence of $Z_1$
with $L^{M^2}$ and use Serre vanishing for its coherent ideal
(Stacks, Lemma 30.16.2, tag \texttt{02O1}).  Choose one linear
right inverse to (A11n) and compose with the prescribed jet
map.  This produces a single linear family of actual algebraic
theta sections realizing every prescribed value and first
derivative.  Consequently it preserves every specified
determinant of first derivatives, including all the finite
linear combinations and class sums in B19 and A10.

The symmetric frame used here has no linear change relative
to the canonical completed frame.  Explicitly, the ratio of
two normalized symmetric frames has value one and is even,
so its differential at zero vanishes.  The quadratic canonical
completion also has no linear term.  Thus the interpolated
first derivatives are the same first derivatives needed for
the canonical finite theta values; higher jets would require
the corresponding higher-order comparison.

After choosing a finite theta structure, let
$J_{\rm new}:W_{\rm new}\to H^0(X,L^{M^2})$ be the algebraic
Heisenberg isomorphism and
\[
        \mathcal T_{\rm new}
               =c_{\rm new}(J_{\rm new}\otimes\C),
        \qquad \dim W_{\rm new}=M^{2\dim X}.
 \tag{A11o}
\]
Applying $J_{\rm new}^{-1}$ to the preceding linear family
gives actual finite theta inputs.  Their analytic first jets
are $c_{\rm new}$ times the prescribed algebraic jets.  The
old analytic jet factors $c_i p_i/N_i$, or the stabilized
factors of A10, remain explicit external comparison scalars.
No assertion that their ratios with $c_{\rm new}$ are
algebraic is needed for this realization.  The right inverse
may have a substantial denominator or norm; it is a recorded
cost and not a uniform quantitative estimate.

There is an integral version at any chosen odd good prime
where the actions, polarization, specified eigenprojections,
and tangent lattice have been extended integrally.  Over a
complete discrete valuation ring, $Z_1$ is finite flat and its
jet module is free.  Relative Serre vanishing makes (A11n)
surjective for sufficiently large $M$ prime to $p$.  A map
onto a finite free module has a linear right inverse, so an
integral prescribed jet family has integral lifts.  Symmetric
references with value one are obtained by averaging with
$[-1]$ and dividing by $2$; their first derivatives vanish.
Prime-to-$p$ level structures can be acquired after a finite
extension.  If an old tangent lattice was not saturated, its
denominator must be retained before applying this assertion.
This argument does not assert an integral identification at a
prime where these hypotheses fail.

\paragraph{What the construction settles.}
This proves an actual generic finite CM host realization,
with a field extension of degree at most four, minimal
rank-one multiplicity padding, an explicit factor-eight
principalization, and a finite theta level realizing the
required first-jet functional.  The added constituents keep
their full CM types.  Equal local multiplicities alone have
not been used to identify two $E$-Hodge structures or two
polarized lattices.  The construction does not identify the
new sections with a prescribed global automorphic family,
preserve its weights or signatures, or supply the separate
definite Gram-one \texttt{StdPair} of the recovered global
chain.  Those comparisons concern more than a finite CM
evaluation at one point.

There is also a precise distinction between two source records
named \texttt{Host}.  A11 constructs
\texttt{ReciprocityTheory.ChartRestriction.G.Geo.Tor.Host}.
The global \texttt{StdPair} is over a \texttt{D7} whose host is
\texttt{Formalism.BasicNotions.Host}, defined in
\texttt{Formalism/LocalRingCardinality/DataHost.lean}.
The latter extends \texttt{HC.HostAx} and additionally
specifies a Galois/Frobenius frame, auxiliary primes and
size conditions, and a factorization of a relative different.
Neither (A11b) nor multiplicity padding is asserted to
construct those additional records.

The construction uses actual homomorphisms, polarizations,
Poincar\'e biextensions, and sections.  In the requested
Grothendieck framework, the correspondence and polarization
viewpoint is the one recorded in \emph{cote44}, scans 7--8;
the dimension and Lefschetz discussion is \emph{cote16}, scans
43--46.  No unproved Chow projector is used here.  The
compatible principalization and finite-jet arguments above
are proved here in full; the explicitly named representability, CM descent,
cube, and cohomology theorems are cited background inputs.
